\clearpage
\section{Exact dyadic capacity and numerical-profile obstruction}\label{app:dyadic_capacity}

Status: ordinary proof from the supplied actual defining sequence and
the retained primary Bushnell--Kutzko results. No Lean or numerical
computation is used. This note calculates the quotient of an actual
simple type. It does not prove universal H5 or identify an explicitly
constructed host with the source-defined arithmetic candidate.

\hypertarget{dyadic_capacity-exact-hypotheses-and-result}{%
\subsection{1. Exact hypotheses and
result}\label{dyadic_capacity-exact-hypotheses-and-result}}

Let F=F\_2((t)). Fix d\textgreater=1 and increasing positive odd
integers

\begin{Verbatim}[fontsize=\small,breaklines=true,breakanywhere=true]
h_1<...<h_d,       q_i=3h_i,       n=2^d.
\end{Verbatim}

Suppose the actual arithmetic constructor supplies a full-degree totally
ramified parameter beta on the complete period-n chain, with actual
approximating fields of degrees

\begin{Verbatim}[fontsize=\small,breaklines=true,breakanywhere=true]
2,4,...,2^d
\end{Verbatim}

and ambient critical levels

\begin{Verbatim}[fontsize=\small,breaklines=true,breakanywhere=true]
nu_j=2^(d−j) q_(d−j+1),       1<=j<=d.
\end{Verbatim}

These are the supplied dyadic specialization of the recursive tower's
field/profile conclusions. The calculation applies equally to the
original polynomial tower of Block166 and to an Artin--Schreier variant
once that variant's field/profile conclusions are established. No
trace-zero assertion is needed here.

The exact actual Heisenberg quotient is

\begin{Verbatim}[fontsize=\small,breaklines=true,breakanywhere=true]
dim_F2(J^1/H^1)=2^d−2,
dim eta=2^((2^d−2)/2).
\end{Verbatim}

Consequently its symplectic dimension is at least the 2d dimensions of
the whole d-plane source packet exactly when d\textgreater=3. In rank~8,
d=3, the quotient has dimension6 and eta has dimension8: the source
carrier fits exactly, with no complementary Heisenberg factor needed.

\begin{longtable}[]{@{}
  >{\raggedright\arraybackslash}p{(\columnwidth - 10\tabcolsep) * \real{0.1667}}
  >{\raggedright\arraybackslash}p{(\columnwidth - 10\tabcolsep) * \real{0.1667}}
  >{\raggedright\arraybackslash}p{(\columnwidth - 10\tabcolsep) * \real{0.1667}}
  >{\raggedright\arraybackslash}p{(\columnwidth - 10\tabcolsep) * \real{0.1667}}
  >{\raggedright\arraybackslash}p{(\columnwidth - 10\tabcolsep) * \real{0.1667}}
  >{\raggedright\arraybackslash}p{(\columnwidth - 10\tabcolsep) * \real{0.1667}}@{}}
\toprule\noalign{}
\begin{minipage}[b]{\linewidth}\raggedright
d
\end{minipage} & \begin{minipage}[b]{\linewidth}\raggedright
Ambient rank
\end{minipage} & \begin{minipage}[b]{\linewidth}\raggedright
Actual quotient dimension
\end{minipage} & \begin{minipage}[b]{\linewidth}\raggedright
eta dimension
\end{minipage} & \begin{minipage}[b]{\linewidth}\raggedright
Source carrier dimension
\end{minipage} & \begin{minipage}[b]{\linewidth}\raggedright
Capacity
\end{minipage} \\
\midrule\noalign{}
\endhead
\bottomrule\noalign{}
\endlastfoot
1 & 2 & 0 & 1 & 2 & Fails \\
2 & 4 & 2 & 2 & 4 & Fails \\
3 & 8 & 6 & 8 & 8 & Exact \\
4 & 16 & 14 & 128 & 16 & Sufficient \\
general & 2\^{}d & 2\^{}d−2 & \(2^{2^{d-1}-1}\) & 2\^{}d &
Sufficient iff d\textgreater=3 \\
\end{longtable}

This improves the leading-layer-only sufficient threshold
d\textgreater=4 to the exact whole-profile threshold d\textgreater=3. It
also exposes a genuine low-rank limitation of this proposed
whole-carrier realization.

\hypertarget{dyadic_capacity-floors-and-ceilings-in-the-actual-bk-orders}{%
\subsection{2. Floors and ceilings in the actual BK
orders}\label{dyadic_capacity-floors-and-ceilings-in-the-actual-bk-orders}}

Write P=rad(A), and let B\_j be the A-centralizer of the degree-2\^{}j
approximation. Its complete-chain grade has F\_2-dimension

\begin{Verbatim}[fontsize=\small,breaklines=true,breakanywhere=true]
b_j=2^(d−j),       b_0=n.
\end{Verbatim}

Put

\begin{Verbatim}[fontsize=\small,breaklines=true,breakanywhere=true]
a_j=ceil(nu_j/2),       t_j=floor(nu_j/2)+1.
\end{Verbatim}

The exact BK definitions, 3.1.7--8, use a\_j as the J cutoff and t\_j as
the H cutoff. Thus a\_j=t\_j for an odd critical level, while
t\_j=a\_j+1 for an even critical level.

The increasing h\_i give

\begin{Verbatim}[fontsize=\small,breaklines=true,breakanywhere=true]
nu_j>2 nu_(j+1),       1<=j<d.
\end{Verbatim}

This strict separation is useful: the centralizer approximation estimate
from BK 3.1.5 puts each newly introduced centralizer into the preceding
one modulo a term above the preceding half-level. In the notation of
these stages, the relevant error level is
nu\_j−nu\_(j+1)\textgreater nu\_j/2. Therefore the successive grade
dimensions really are b\_j; there is no unproved assumption that the
approximation fields are literally nested inside the final field.

Applying the ring recursion gives the J-grade staircase:

\begin{itemize}
\tightlist
\item
  below a\_d the grade dimension is b\_d=1;
\item
  on a\_(j+1)\textless=k\textless a\_j it is b\_j;
\item
  from a\_1 upward it is b\_0=n.
\end{itemize}

The H staircase uses the corresponding t\_j. Hence the exact
contribution of transition j to \(J^1/H^1\) is

\begin{Verbatim}[fontsize=\small,breaklines=true,breakanywhere=true]
(b_(j−1)−b_j) · 1_(nu_j even).
\end{Verbatim}

One can also obtain this by subtracting the exact logarithmic indices of
H and J in 1+P. It is an equality of the group index
\textbar \(J^1/H^1\)\textbar, not a guess from the visible
leading matrix support. BK 3.4.1 makes the resulting commutator pairing
nondegenerate.

For j\textless d the level nu\_j has a factor of2 and is even. The final
level nu\_d=q\_1 is odd. Therefore

\begin{Verbatim}[fontsize=\small,breaklines=true,breakanywhere=true]
dim_F2 J^1/H^1
  =sum_(j=1)^(d−1) (2^(d−j+1)−2^(d−j))
  =2^d−2.
\end{Verbatim}

The omitted final increment would have dimension1. It cannot be treated
as a symplectic plane or silently restored: at that odd level the two
defining lattice cutoffs are exactly equal. Its disappearance is also
consistent with the even dimension required of a nondegenerate
alternating space.

The source packet requires dimension2d. At d=1,2, 2\^{}d−2\textless2d.
At d=3 equality holds. If the inequality holds for d\textgreater=3, then

\begin{Verbatim}[fontsize=\small,breaklines=true,breakanywhere=true]
2^(d+1)−2−2(d+1) = (2^d−2−2d)+(2^d−2)>0,
\end{Verbatim}

so it continues to hold for every larger d.~This proves the exact
threshold.

\hypertarget{dyadic_capacity-conductors-still-match-despite-the-lost-final-contribution}{%
\subsection{3. Conductors still match despite the lost final
contribution}\label{dyadic_capacity-conductors-still-match-despite-the-lost-final-contribution}}

This section reuses the same retained formal-degree and self-pair
conductor identities already audited in Blocks161 and165. It is not
needed to prove the dimension formula, but shows why numerical conductor
checks do not expose the low-rank capacity failure.

Put Delta\_j=b\_(j−1)−b\_j. The exact grade count gives

\begin{Verbatim}[fontsize=\small,breaklines=true,breakanywhere=true]
I=log_2[(1+P):J^1]
  =sum_j Delta_j ceil(nu_j/2)−(n−1).
\end{Verbatim}

For this profile only the final level is odd and Delta\_d=1. Therefore

\begin{Verbatim}[fontsize=\small,breaklines=true,breakanywhere=true]
2I=sum_j Delta_j nu_j+1−2(n−1).
\end{Verbatim}

The Heisenberg dimension exponent is D/2, where D=n−2. The retained
formal-degree identity for a full-degree, totally ramified, torsion-one
maximal type gives

\begin{Verbatim}[fontsize=\small,breaklines=true,breakanywhere=true]
a(Pi^vee×Pi)=n^2−1+(n−1)+2I+D.
\end{Verbatim}

Substitution yields exact cancellation of the odd-cutoff correction:

\begin{Verbatim}[fontsize=\small,breaklines=true,breakanywhere=true]
a(Pi^vee×Pi)=n^2−1+sum_j Delta_j nu_j
           =2^(2d)−1+sum_(i=1)^d 3·4^(i−1) h_i.
\end{Verbatim}

The ordinary conductor is n+nu\_1, so Swan(Pi)=2\^{}(d−1)·3h\_d, and the
torsion number is one. These are the usual three numerical invariants of
the stipulated independent dyadic Heisenberg tensor source.

Thus all three numerical invariants can match while eta is too small for
the whole source carrier at ranks2 and4. This is a capacity failure, not
a conductor mismatch. It does not by itself refute any original H1--H5
statement.

\hypertarget{dyadic_capacity-rank~2-failure-is-forced-by-swan-and-torsion-alone}{%
\subsection{4. Rank2 failure is forced by Swan and torsion
alone}\label{dyadic_capacity-rank~2-failure-is-forced-by-swan-and-torsion-alone}}

Let Pi be any actual supercuspidal representation of GL2(F) with torsion
number one and odd Swan exponent N=3h\_1. BK 6.2.5 forces its maximal
parameter field K/F to be totally ramified of degree~2 and its hereditary
order to have complete period2. The ordinary conductor formula then
fixes its native level as N.

Every parameter beta of a maximal simple type in Pi has v\_K(beta)=−N,
which is odd. Since the residue extension is trivial, BK 1.4.14--15 say
beta is minimal over F. The two order definitions are therefore

\begin{Verbatim}[fontsize=\small,breaklines=true,breakanywhere=true]
J_order=O_K+P^ceil(N/2),
H_order=O_K+P^(floor(N/2)+1).
\end{Verbatim}

The exponents coincide because N is odd. Hence \(J^1=H^1\),
eta is one-dimensional, and no two-dimensional symplectic source plane
can embed into this quotient.

This proves failure for every maximal simple type of every candidate
with these two invariants. It is stronger than failure of one explicit
matrix formula. It remains a statement about attachment inside the
actual maximal-simple-type Heisenberg carrier, not a proof that no other
arithmetic realization relation could work.

\hypertarget{dyadic_capacity-rank~4-failure-is-also-forced-by-swan-and-torsion}{%
\subsection{5. Rank4 failure is also forced by Swan and
torsion}\label{dyadic_capacity-rank~4-failure-is-also-forced-by-swan-and-torsion}}

Let Pi be any actual supercuspidal representation of GL4(F) with torsion
number one and Swan exponent N=6h\_2, where h\_2 is odd. Its maximal
parameter field is totally ramified of degree~4, its order has period4,
and its native level is N.

Here v\_2(N)=1. The defining-sequence/minimality argument of Block166
forces the terminal minimal approximation field E/F to have degree~2.
There is no possible intermediate degree between2 and4. The scalar
terminal case is excluded by 4 not dividing N, and a minimal
full-degree~4 parameter is excluded because gcd(N,4)=2. Thus the actual
degree pattern is forced to be 4→2.

Write c=−k\_0(beta,A), so 1\textless=c\textless N. Select the critical
approximation furnished by the construction in BK 2.4.1 and the remark
preceding 2.4.7. Because the original field is maximal, this
approximation can be selected with a simple maximal derived stratum in
B=End\_E(F\^{}4)≅M2(E). Equivalently one can replace its derived stratum
by an equivalent simple representative and use BK 2.2.2--3. The derived
field has degree~2 and ramification index2 over E. A simple stratum
{[}B\_order,c,c−1,alpha{]} is minimal. Since the complete B-chain has
period2 and the derived field is totally ramified, BK 1.4.14 forces

\begin{Verbatim}[fontsize=\small,breaklines=true,breakanywhere=true]
gcd(c,2)=1.
\end{Verbatim}

Therefore c is odd. It cannot be made even while preserving this degree
pattern and a genuine maximal derived field.

Let j=(c+1)/2. The actual orders are

\begin{Verbatim}[fontsize=\small,breaklines=true,breakanywhere=true]
J_order=O_K+Q^j+P^(N/2),
H_order=O_K+Q^j+P^(N/2+1).
\end{Verbatim}

The final low transition contributes nothing, while the high transition
has dimension4−2=2. Thus every such Pi has

\begin{Verbatim}[fontsize=\small,breaklines=true,breakanywhere=true]
dim_F2 J^1/H^1=2,       dim eta=2.
\end{Verbatim}

The desired two-plane source carrier has dimension4, so it cannot be a
sector of this eta. Requiring the self-pair conductor additionally fixes
c=3h\_1 through

\begin{Verbatim}[fontsize=\small,breaklines=true,breakanywhere=true]
a(Pi^vee×Pi)=15+2N+c,
\end{Verbatim}

but is unnecessary for the capacity failure.

\hypertarget{dyadic_capacity-rank~8-the-three-invariants-do-not-force-the-good-profile}{%
\subsection{6. Rank8: the three invariants do not force the good
profile}\label{dyadic_capacity-rank~8-the-three-invariants-do-not-force-the-good-profile}}

The full 2→4→8 tower has quotient dimension6 and enough capacity.
Nevertheless Swan, torsion number and self-pair conductor alone do not
force that degree pattern. Here is an actual skipped-degree arithmetic
construction with the same three numerical invariants and a smaller
quotient.

Fix increasing odd h\_1\textless h\_2\textless h\_3 and put

\begin{Verbatim}[fontsize=\small,breaklines=true,breakanywhere=true]
q=3h_3,       N=4q=12h_3,       c=h_1+4h_2.
\end{Verbatim}

Then c is odd and 0\textless c\textless N; indeed c\textless5h\_3, so
even 2c\textless N.

Construct the separable totally ramified quadratic extension

\begin{Verbatim}[fontsize=\small,breaklines=true,breakanywhere=true]
E=F(alpha),       alpha^2+alpha=t^(−q).
\end{Verbatim}

Put a=(q+1)/2 and z=t\^{}a alpha. Then z is a uniformizer and

\begin{Verbatim}[fontsize=\small,breaklines=true,breakanywhere=true]
z^2+t^a z+t=0.
\end{Verbatim}

The element alpha has E-valuation −q and is minimal over F. On E\^{}4,
gamma=alpha I\_4 is minimal of ambient level N in the complete
period-eight order A.

Now construct the separable Eisenstein quartic extension over E

\begin{Verbatim}[fontsize=\small,breaklines=true,breakanywhere=true]
L=E(e),       e^4+z e+z=0,
\end{Verbatim}

and let C be multiplication by e\^{}(−c) on L as a four-dimensional
E-space. Since c is odd, e\^{}(−c) is minimal, generates L/E, and
defines a simple maximal derived stratum {[}B\_order,c,c−1,C{]}.

Let P0 extract the coefficient of1 on each E-coordinate in the basis
1,z. The normalized outer corestriction from the Artin--Schreier
right-inverse calculation has

\begin{Verbatim}[fontsize=\small,breaklines=true,breakanywhere=true]
s(P0)=t/z^2.
\end{Verbatim}

Consequently the explicit F-linear lift

\begin{Verbatim}[fontsize=\small,breaklines=true,breakanywhere=true]
b=(z^2/t) C P0
\end{Verbatim}

satisfies s(b)=C and has valuation −c.~Define the actual rank-eight
matrix

\begin{Verbatim}[fontsize=\small,breaklines=true,breakanywhere=true]
beta=alpha I_4+(z^2/t) C P0.
\end{Verbatim}

BK 2.2.3 applies because gamma is simple at level c\textless N and the
derived stratum is simple maximal. It proves that F{[}beta{]}/F is
totally ramified of degree~8, normalizes the complete chain, has leading
level N, and has critical exponent −c.~Its defining degree pattern is
8→2, skipping degree~4.

Since N is even and c is odd, the actual Heisenberg quotient has
dimension

\begin{Verbatim}[fontsize=\small,breaklines=true,breakanywhere=true]
(8−4)·1+(4−1)·0=4,
\end{Verbatim}

so eta has dimension4. Yet its invariants are

\begin{Verbatim}[fontsize=\small,breaklines=true,breakanywhere=true]
Swan=N=12h_3,
torsion=1,
a(Pi^vee×Pi)=63+4N+3c
             =63+48h_3+12h_2+3h_1,
\end{Verbatim}

exactly the same as the full three-stage profile. Any maximal simple
character and its standard type completion give an actual supercuspidal
Pi with these values. The required full three-plane source carrier would
have dimension8, larger than this eta.

This is an arithmetic counterexample to the proposed implication from
those three numerical invariants to adequate rank-eight Heisenberg
capacity. It is not an H5 counterexample, does not identify the skipped
host with the source-defined classical candidate, and does not claim
that it preserves the supplied full source characters.

\hypertarget{dyadic_capacity-which-repairs-have-actually-been-tested}{%
\subsection{7. Which repairs have actually been
tested}\label{dyadic_capacity-which-repairs-have-actually-been-tested}}

\hypertarget{dyadic_capacity-replacing-the-displayed-type-or-parameter-at-fixed-swan-and-torsion}{%
\subsubsection{Replacing the displayed type or parameter at fixed Swan
and
torsion}\label{dyadic_capacity-replacing-the-displayed-type-or-parameter-at-fixed-swan-and-torsion}}

At ranks2 and4 this fails by Sections~4--5 for every maximal simple type
satisfying those invariants. A new matrix search or a different critical
coefficient cannot manufacture the missing quotient dimensions.

At rank~8 the full tower supplies the good profile, but the three
invariants alone do not force it inside a separately selected candidate.
Section~6 proves the failure of that proposed inference with an explicit
actual host. A finer endo-class or ramification comparison would be
needed for that same-candidate assertion.

\hypertarget{dyadic_capacity-skipping-approximation-degrees-or-changing-critical-parity}{%
\subsubsection{Skipping approximation degrees or changing critical
parity}\label{dyadic_capacity-skipping-approximation-degrees-or-changing-critical-parity}}

At rank~2 there is no degree to skip. At rank~4 the level excludes both a
scalar minimal approximation and a full-degree minimal approximation.
The remaining quadratic derived step has odd critical exponent by
minimality. Thus neither skipping nor declaring that exponent even is a
valid repair at fixed invariants.

At rank~8 a degree skip is genuinely possible, but it decreases capacity
in the explicit example above.

\hypertarget{dyadic_capacity-scalar-twists-shallow-or-deep}{%
\subsubsection{Scalar twists, shallow or
deep}\label{dyadic_capacity-scalar-twists-shallow-or-deep}}

A scalar twist of the existing actual inducing representation does not
increase its dimension. The exact Block164 pushout transport identifies
the same quotient J/H and leaves its commutator pairing unchanged. Thus
it cannot turn an insufficient eta into a whole source carrier.

For an actual parameter, replacing beta by beta+aI with a∈F leaves its
field, centralizer and commutator map unchanged. A shallow scalar
addition does not change the leading level. A dominating scalar addition
can change that level, but it contributes no new noncentral Heisenberg
directions: in the defining sequence the added scalar stage has equal H
and J orders. This agrees with the exact representation-level twist
transport. Enlarging a displayed numerical level is therefore not a
capacity repair.

This paragraph does not assert that every arbitrary scalar character
twist is represented by one globally specified additive parameter. The
already proved group/representation transport suffices for its negative
capacity conclusion.

\hypertarget{dyadic_capacity-adding-residue-multiplicity-or-increasing-ambient-rank}{%
\subsubsection{Adding residue multiplicity or increasing ambient
rank}\label{dyadic_capacity-adding-residue-multiplicity-or-increasing-ambient-rank}}

The torsion-one condition forces a full-degree totally ramified
parameter field and residue multiplicity one. A host with an unramified
or residual multiplicity factor does not preserve that invariant.
Increasing the ambient rank changes the required source rank. Neither
maneuver closes the current fixed-rank construction.

\hypertarget{dyadic_capacity-using-a-larger-arithmetic-carrier-inside-the-same-representation}{%
\subsubsection{Using a larger arithmetic carrier inside the same
representation}\label{dyadic_capacity-using-a-larger-arithmetic-carrier-inside-the-same-representation}}

A larger compact subgroup representation, a genuinely different
arithmetic quotient, or an induction-stage carrier might have enough
dimension. Such an enlargement is not ruled out here. It would need an
actual finite Heisenberg action carrying the prescribed full source data
and a proved recovery in the unchanged candidate. Artificially adding
vector-space copies or transporting abstract operators to an induced
space does not provide that arithmetic attachment or prove its kernel
compatibility. No such repair is constructed in this note.

\hypertarget{dyadic_capacity-quadratic-refinements-and-marked-centers}{%
\subsection{8. Quadratic refinements and marked
centers}\label{dyadic_capacity-quadratic-refinements-and-marked-centers}}

For d\textgreater=3, the dimension result gives a symplectic embedding
of the source quotient into the actual host quotient. With C trivial, a
C×-central pushout can then be matched by its alternating commutator,
and Stone--von Neumann supplies the corresponding source sector. This is
the same finite scalar-pushout attachment contract used previously.

At p=2, that statement must not be promoted to an identification of
arbitrary preselected native quadratic refinements. The C× pushout
permits rescaling lifts by i and can identify the pushouts of the
dihedral and quaternion extraspecial planes, even though their native
square functions differ. A chosen source embedding can preserve that
source's full finite central extension and its character values. It does
not thereby identify a separately marked target mu\_2 extension or its
Arf invariant.

In particular, symplectic-isotropic reduction respects a prescribed
quadratic form only when the reduced subspace is totally singular for
that form. The retained \path{DYADIC_AND_WEAK_ATTACHMENT159.md} gives the
exact quadratic/Arf obstruction. Additional source or target marked
centers, chosen lifts, quadratic refinements and nontrivial stabilizer
actions therefore require their own comparison. They are neither erased
from the task nor supplied by a dimension inequality.

\hypertarget{dyadic_capacity-scope-and-source-provenance}{%
\subsection{9. Scope and source
provenance}\label{dyadic_capacity-scope-and-source-provenance}}

The strongest proved conclusion is exact capacity for the specified
actual full defining profile, together with invariant low-rank failures
and an explicit rank-eight counterexample to numerical profile forcing.
No universal impossibility for original Theorem~16.2 is asserted. In
particular this note supplies no complete original H1--H5 counterexample
and does not decide whether that endpoint requires this precise
whole-maximal-type-carrier realization.

Primary source directly reread: Bushnell--Kutzko, \emph{The admissible
dual of GL(N) via compact open subgroups}, 1.4.14--15; 2.2.2--3;
2.4.1--2 and the maximal-field remark immediately preceding 2.4.7;
3.1.5, 3.1.7--9; 3.2.1--4; 3.4.1; 6.2.5 and its following paragraph.

Authorized retained text:

\begin{Verbatim}[fontsize=\small,breaklines=true,breakanywhere=true]
closure_recovery/extracted/Paper_XI_Continuation158_General_Sectors_and_Arit
hmetic_Constraints/scratch158/information158/input_cache/BK_source.txt
SHA256 d69db0827256707a34cd2562af68b9de11fc889ab3f43535ec47b620d58ca9c8
\end{Verbatim}

Internal inputs reused at their stated scope:

\begin{itemize}
\tightlist
\item
  \path{attachment166/tower/RECURSIVE_TOWER166.md}: actual defining degrees,
  levels, filtered right inverse and separated grade count; its field
  constructor permits p=2.
\item
  \path{attachment166/tower/TRACE_ZERO_AS_TOWER166.md}, Sections~1--2:
  normalized corestriction and explicit right inverse. Their algebraic
  formulas apply at p=2; the odd-prime trace-zero assertion is not used.
\item
  \path{attachment166/capacity/LEADING_LAYER_CAPACITY166.md}: the leading
  subspace argument and the minimal terminal degree calculation.
\item
  \path{base_realization165/polarization/PROFILE_RIGIDITY165.md} and
  \path{closure_recovery/full161/scratch161/placement161/ADJOINT_CONDUCTOR_TEST161.md}:
  retained formal-degree/self-pair conductor identities and
  normalization.
\item
  \path{input_pair164/TWIST_TRANSPORT.md}: exact character-twist/pushout
  transport.
\item
  \path{closure_recovery/extracted/Paper_XI_Continuation159_Type_Completion_and_Sector_Recovery/scratch159/coverage159/DYADIC_AND_WEAK_ATTACHMENT159.md}:
  distinction between symplectic and marked quadratic reduction.
\end{itemize}

No full copyrighted source is redistributed here. No new external
correspondence theorem or H5 premise is introduced.
